\section{Discussion}

\subsection{What a practitioner should take from this}

The first decision in an evaluation is which convention to score under, and on
this data that decision moves the answer more than the choice of adaptation
policy does. It should therefore be made deliberately and stated. Inheriting it
from whichever paper the code was adapted from is how it usually happens.

The convention follows from what an alarm costs. If an operator investigates
each alarm individually, every false row is paid for and point-wise precision is
the quantity that matters. If an operator investigates each \emph{episode}, a
detector that flags one row of a fault has done its job and the composite or
range-based view describes the deployment better. These are different questions
about the same predictions. Reporting one number without saying which question
it answers leaves the reader to guess.

The second suggestion is procedural. Quote a trivial baseline next to every
score. At a base rate of \BaseRate\%, flagging every row scores \FlagAll{} under
point-wise F1, which is higher than several published-looking numbers. We spent
part of this project rediscovering that a detector can clear 0.5 F1 while doing
nothing at all, and the baseline column is what stops that from happening
quietly.

The third point is about which number to read once both of those are settled.
The strategy that wins our headline table on the mean is not the one we would
deploy. \BestComposite{} takes the top place on average and still fails to beat
the trivial baseline on \BestCompositeLosses{} of the \NRecordings{} individual
recordings, with an interquartile range that straddles zero;
\SteadiestName{} scores lower on average and fails on \SteadiestLosses{}. A mean
over streams answers a question about a benchmark. An operator with one stream
is asking a different question, and the spread across recordings is what answers
it.

\subsection{On the adaptation question itself}

Stripped of the metric problem, the practical picture is narrower than
``adaptive detection works better''. Rebuilding on a schedule and rebuilding on
detected change both beat leaving the model alone, and the margin between the
two is small next to the margin between either of them and the trivial
baseline. The case for drift-triggered rebuilding is not that it scores higher.
It is that it reacts to evidence and not to a clock, which lets it notice a
sudden shift within a few rows instead of waiting for the next scheduled
rebuild, and keeps it from rebuilding on a stream that has not moved.

That argument is about the shape of the response, and Section~\ref{sec:timing}
is where it is measured. It is a smaller claim than the one we expected to make
when the project started.

\subsection{Limitations}
\label{sec:limitations}

These are constraints on what the numbers above can support, and several of them
cut against our own results.

\paragraph{The base recording is not stationary.} SKAB's anomaly-free recording
has no annotated faults but the process drifts on its own. Running with no
injected drift still produces adaptations, at nearly the same positions as with a
3\,sd step injected. False-alarm counts on the injected streams are therefore an
upper bound and not a measurement: some of those firings are the detector
correctly noticing real process change for which we have no ground truth. For
the same reason we do not report detection \emph{rate}, which comes out at 1.0
for every magnitude including the zero control, because a firing always lands
within the matching horizon of an arbitrary point.

\paragraph{The labelled stream is concatenated.} The individual \texttt{valve1}
recordings are about 1{,}000 rows, too short to warm up on and still leave stream
to measure, so we join 16 of them. The joins are genuine regime
changes of unknown magnitude. That is why this stream is used for anomaly
quality only and never for drift timing, and it is why we also report per
recording, where the joins are not inside the scored region.

\paragraph{Event recall saturates on SKAB.} SKAB's faults are long, so on the
labelled stream every strategy catches every event and event-level recall is 1.0
across the board. Composite F1 there reduces to a precision-driven ranking. That
is the correct operational picture, and it is what point-wise scoring obscured,
but the recall axis carries no information on that particular stream. This is
why the UCR subset is in the paper, and Section~\ref{sec:ucr} is where the
recall axis does some work. The limitation is not removed so much as bounded: it
applies to SKAB, partly to NAB, and not to UCR.

\paragraph{Three datasets, one of them contested.} \citet{wu2023flawed} argue
that NAB and benchmarks like it are flawed enough to mislead. We use NAB only to
ask whether a calibration transfers, which is a weaker use than ranking
detectors on it, but the criticism still applies to Section~\ref{sec:transfer}.
The UCR archive is those authors' own replacement and we use a subset of it, so
the study is no longer resting on the benchmarks they criticised. It is a
subset, though: \NUcrSeries{} of 250, chosen by a stated rule and bounded by
what one core can run. A study with the compute to take all 250 would be a
stronger test than ours.

\paragraph{No deep or model-pooling baselines.} The comparison covers one online
model, one static ensemble and three classical detectors. It does not include the
autoencoder and model-pooling systems that dominate recent leaderboards
\citep{yoon2022arcus}. Our claim is about how adaptation policies compare under
different scoring conventions, not about the state of the art in raw detection
quality, and the metric result is if anything more consequential for those
systems, since they are usually reported under point-adjust.

\paragraph{The threshold is a deployment choice.} It has to be set from the
expected anomaly rate of the specific stream. We measured what carrying one
setting across datasets costs instead of assuming it away, but no rule here
removes the need to choose.

\paragraph{One machine, one implementation.} Throughput in
Section~\ref{sec:cost} is from a single laptop core and a specific pair of
libraries. It supports statements about relative cost between the strategies and
not about absolute deployability.

\subsection{Reproducing this}

Every number in this paper, including the ones in the prose, is generated from
the committed CSV files by \texttt{scripts/make\_paper\_tables.py}; the
manuscript contains no hand-typed values. The experiments regenerate with

\begin{quote}
\texttt{python scripts/run\_multiseed.py}\\
\texttt{python scripts/run\_sweep.py -{}-seeds 5}\\
\texttt{python scripts/run\_ablation.py}\\
\texttt{python scripts/run\_throughput.py}\\
\texttt{python scripts/make\_paper\_figures.py}\\
\texttt{python scripts/make\_paper\_tables.py}
\end{quote}

Timing gets its own short run because it is the one measurement sensitive to
what else the machine is doing, and the grid occupies a core for the better part
of an hour.

The full grid takes under an hour on one core. The test suite covers the metric
implementations against hand-computed values, the statistics against constructed
data with known answers, and the fast scoring path against the plain one.

\section{Conclusion}

We compared four adaptation policies for a streaming anomaly detector, crossed
with the model they run on, repeated over \NSeeds{} seeds and \NRecordings{}
recordings, and scored the same runs under four published conventions. The four
conventions produce \NDistinctWinners{} different winners.

The narrow finding is about one benchmark and one family of detectors. The
general one is that a paper reporting a single scoring convention has not
established which policy is better, only which policy that convention prefers.
On the stream studied here, \MoverName{} moves from \MoverPointRank{} to
\MoverCompositeRank{} of \NStrategies{} on identical predictions.

The third dataset put a limit on all of it. On the UCR subset no policy
separates from the others, the classical static detectors match anything with a
rebuild schedule, and under the archive's own criterion the field as a whole
puts its peak score inside the labelled window on \MeanUcrTopOne{} of series.
The comparison this paper is built around is worth making on SKAB and is close
to meaningless on the benchmark built to be hard. We would rather finish on that
than on the part that worked.
